\chapter{Metoda i implementacja systemu}
\label{ch:metoda-implementacja}

\section{Założenia projektowe proponowanego rozwiązania}
\label{sec:zalozenia-projektowe}

% TODO: Opisać wymagania projektowe kodeka: przepływność poniżej 1 kbps, sygnał mowy 16 kHz,
% możliwość pracy w warunkach brzegowych, ograniczony koszt obliczeniowy, rekonstrukcja z dyskretnych indeksów
% oraz rozdzielenie etapu uczenia od etapu inferencji.

\section{Dobór reprezentacji sygnału wejściowego}
\label{sec:reprezentacja-wejsciowa}

% TODO: Uzasadnić pracę bezpośrednio na przebiegu czasowym oraz wykorzystanie reprezentacji STFT i mel
% jako składników funkcji kosztu. Należy wskazać parametry segmentów, próbkowania i ewentualnego przygotowania danych.

\section{Architektura kodera i dekodera}
\label{sec:architektura-kodera-dekodera}

% TODO: Opisać moduły SirenCodec: konwolucyjny koder z redukcją czasu, przestrzeń latentną,
% residual vector quantization, dekoder z nadpróbkowaniem, aktywacje SnakeBeta oraz opcjonalne bloki
% samouwagi, temporalne, pasmowe, harmoniczne i post-refiner.

\section{Proces kompresji i rekonstrukcji sygnału}
\label{sec:kompresja-rekonstrukcja}

% TODO: Przedstawić przepływ danych: audio -> koder -> wektory latentne -> RVQ -> strumień indeksów
% -> dekoder -> sygnał zrekonstruowany. Wzór na nominalną przepływność powinien wynikać z liczby słowników,
% rozmiarów codebooków i częstotliwości ramek latentnych.

\section{Implementacja systemu oraz wykorzystane technologie}
\label{sec:implementacja-technologie}

% TODO: Opisać strukturę repozytorium, backend CUDA/PyTorch, pomocniczy backend MLX, konfiguracje JSON,
% skrypty treningowe, Docker, zapisywanie checkpointów, logów, spektrogramów i plików audio oraz testy automatyczne.
